\documentclass[11pt]{article}
\usepackage{mathblatt}


\blattfuss{Lineare}{Lösungen}
\begin{document}
\noindent{\large\bfseries Lineare \textperiodcentered{} Lösungen}\par\medskip
\begin{pfloesung}{}
\lz{1.}{$A(3|2)$}{}{}
\end{pfloesung}
\begin{pfloesung}{}
\lz{2.}{$B$: $4$ nach rechts, $1$ nach oben}{}{}
\end{pfloesung}
\begin{pfloesung}{}
\lz{3.}{$9$}{g(4) = 8 + 1}{}
\end{pfloesung}
\begin{pfloesung}{}
\lz{4.}{Gerade durch $(0|2)$ und $(2|4)$}{}{}
\end{pfloesung}
\begin{pfloesung}{}
\lz{5.}{$x^2 + 8x + 16$}{}{}
\end{pfloesung}
\begin{pfloesung}{}
\lz{6.}{$x^2 + 18x + 81$}{}{}
\end{pfloesung}
\begin{pfloesung}{}
\lz{7.}{$4$}{x}{}
\end{pfloesung}
\begin{pfloesung}{}
\lz{8.}{$12$}{x}{}
\end{pfloesung}
\begin{pfloesung}{}
\lz{9.}{$S(2|3)$}{x + 1 $\Rightarrow$ -x + 5, x = 2}{}
\end{pfloesung}
\begin{pfloesung}{}
\lz{10.}{$S(3|6)$}{2x = x + 3, x = 3}{}
\end{pfloesung}
\begin{pfloesung}{}
\lz{11.}{$8$}{}{}
\end{pfloesung}
\begin{pfloesung}{}
\lz{12.}{$11$}{}{}
\end{pfloesung}
\begin{pfloesung}{}
\lz{13.}{$x_1 = 6$, $x_2 = -6$}{}{}
\end{pfloesung}
\begin{pfloesung}{}
\lz{14.}{$x_1 = 10$, $x_2 = -10$}{}{}
\end{pfloesung}
\begin{pfloesung}{}
\lz{15.}{$-5$}{3, n}{}
\end{pfloesung}
\begin{pfloesung}{}
\lz{16.}{steigt}{m ist positiv}{}
\end{pfloesung}
\begin{pfloesung}{}
\lz{17.}{$3$}{3 nach oben, also m}{}
\end{pfloesung}
\begin{pfloesung}{{\small\color{mbgrau}P10 2019 · 2b}}
\lz{18.}{Anstieg $-$2: f; parallel zur x-Achse: h}{Steigungsdreieck an f: 1 nach rechts, 2 nach unten}{}
\end{pfloesung}
\begin{pfloesung}{}
\lz{19.}{$3$}{n}{}
\end{pfloesung}
\begin{pfloesung}{}
\lz{20.}{$2$}{m; 1 nach rechts, 2 nach oben}{}
\end{pfloesung}
\begin{pfloesung}{}
\lz{21.}{$-2$}{m; 1 nach rechts, 2 nach unten}{}
\end{pfloesung}
\begin{pfloesung}{}
\lz{22.}{$3x - 4$}{f(x)}{}
\end{pfloesung}
\begin{pfloesung}{{\small\color{mbgrau}P10 2016 · 1c}}
\lz{23.}{f}{}{}
\end{pfloesung}
\begin{pfloesung}{{\small\color{mbgrau}P10 2024 · 1i}}
\lz{24.}{Punkt (0|$-$1) markiert}{$-$1}{}
\end{pfloesung}
\begin{pfloesung}{{\small\color{mbgrau}P10 2025 · 1f}}
\lz{25.}{erste Abbildung (fallende Gerade)}{Gerade ohne Knick}{}
\end{pfloesung}
\begin{pfloesung}{}
\lz{26.}{$\tfrac{1}{2}$}{m; 2 nach rechts, 1 nach oben}{}
\end{pfloesung}
\begin{pfloesung}{{\small\color{mbgrau}P10 2019 · 2a}}
\lz{27.}{x = $-$2}{}{}
\end{pfloesung}
\begin{pfloesung}{{\small\color{mbgrau}P10 2019 · 2c}}
\lz{28.}{f: y = $-$2x + 2; g: y = x + 2; h: y = $-$2}{y-Achsenabschnitt 2 bei f und g}{}
\end{pfloesung}
\begin{pfloesung}{}
\lz{29.}{$-5$, $-2$, $1$, $4$}{}{}
\end{pfloesung}
\begin{pfloesung}{}
\lz{30.}{Start $(0|4)$, 1 nach rechts, 3 nach unten: Gerade durch $(0|4)$ und $(1|1)$}{}{}
\end{pfloesung}
\begin{pfloesung}{}
\lz{31.}{Start $(0|3)$, 2 nach rechts, 1 nach oben: Gerade durch $(0|3)$ und $(2|4)$}{}{}
\end{pfloesung}
\begin{pfloesung}{}
\lz{32.}{Start $(0|-4)$, 1 nach rechts, 2 nach oben: Gerade durch $(0|-4)$ und $(1|-2)$}{}{}
\end{pfloesung}
\begin{pfloesung}{{\small\color{mbgrau}P10 2021 · 2a}}
\lz{33.}{Gerade durch (0; 1) und (1; 4)}{y-Achsenabschnitt 1, Steigungsdreieck 1 nach rechts, 3 nach oben}{}
\end{pfloesung}
\begin{pfloesung}{{\small\color{mbgrau}P10 2024 · 3a}}
\lz{34.}{Gerade durch (0|1) und (1|5), steil steigend}{Steigungsdreieck: 1 nach rechts, 4 nach oben}{}
\end{pfloesung}
\begin{pfloesung}{{\small\color{mbgrau}P10 2022 · 3a}}
\lz{35.}{Gerade durch (0|3) und (1|1); x = 1,5}{0 = $-$2x + 3}{}
\end{pfloesung}
\begin{pfloesung}{}
\lz{36.}{$3x + 4$}{(10 - 4) : (2 - 0) $\Rightarrow$ 3, n = 4: f(x)}{}
\end{pfloesung}
\begin{pfloesung}{}
\lz{37.}{$2x - 1$}{(9 - 3) : (5 - 2) $\Rightarrow$ 2, n = -1: f(x)}{}
\end{pfloesung}
\begin{pfloesung}{}
\lz{38.}{$-3x + 9$}{(0 - 6) : (3 - 1) $\Rightarrow$ -3, n = 9: f(x)}{}
\end{pfloesung}
\begin{pfloesung}{}
\lz{39.}{Gerade durch $A(-3|-1)$ und $B(2|4)$}{}{}
\end{pfloesung}
\begin{pfloesung}{}
\lz{40.}{$2x - 3$}{2, n = -3: f(x)}{}
\end{pfloesung}
\begin{pfloesung}{{\small\color{mbgrau}P10 2017 · 5a}}
\lz{41.}{Gerade durch K und L; m = ½, n = 1 $\Rightarrow$ y = ½x + 1}{(2 $-$ ($-$1)) : (2 $-$ ($-$4)) $\Rightarrow$ 0{,}5; 2 = ½ · 2 + n = 1}{}
\end{pfloesung}
\begin{pfloesung}{{\small\color{mbgrau}P10 2015 · 4d}}
\lz{42.}{h(t) = $-$5t + 1300}{(700 $-$ 1200) : (120 $-$ 20) $\Rightarrow$ $-$5; 1200 = $-$5 · 20 + n = 1300}{}
\end{pfloesung}
\begin{pfloesung}{}
\lz{43.}{$4{,}20$: Grundgebühr 4{,}20 €}{(25{,}20 - 12{,}60) : (10 - 4) $\Rightarrow$ 2{,}10; Unterschied: $12{,}60 - 4 \cdot 2{,}10$ $\Rightarrow$ 4{,}2}{}
\end{pfloesung}
\begin{pfloesung}{{\small\color{mbgrau}P10 2025 · 5a}}
\lz{44.}{Gerade durch A($-$2|6) und B(3|$-$1,5); Aussage 1 falsch (fallend), Aussage 2 wahr; f(x) = $-$1,5x + 3}{($-$1,5 $-$ 6) : (3 + 2) $\Rightarrow$ $-$1,5; 3}{}
\end{pfloesung}
\begin{pfloesung}{}
\lz{45.}{$10$}{Unterschied: $3 \cdot 4 - 2$ $\Rightarrow$ 10}{}
\end{pfloesung}
\begin{pfloesung}{}
\lz{46.}{Gerade}{}{}
\end{pfloesung}
\begin{pfloesung}{}
\lz{47.}{Normalform}{}{}
\end{pfloesung}
\begin{pfloesung}{}
\lz{48.}{$x + 4$}{$x^2 - 2x$}{}
\end{pfloesung}
\begin{pfloesung}{}
\lz{49.}{$13$}{f(4) = $2 \cdot 4 + 5$}{}
\end{pfloesung}
\begin{pfloesung}{}
\lz{50.}{Schnittpunkt $(0|5)$}{f(0) = 5}{}
\end{pfloesung}
\begin{pfloesung}{}
\lz{51.}{$x_1 = 3$, $x_2 = -1$}{0 = $(x - 1)^2 - 4, 4$ = $(x - 1)^2, x - 1$ = 2 oder x - 1 = -2}{}
\end{pfloesung}
\begin{pfloesung}{}
\lz{52.}{$36$}{f(-6) = $(-6)^2$}{}
\end{pfloesung}
\begin{pfloesung}{}
\lz{53.}{$-5$}{f(-3) = 9 - 12 - 2}{}
\end{pfloesung}
\begin{pfloesung}{{\small\color{mbgrau}P10 2017 · 5b}}
\lz{54.}{y = $-$4}{½ · ($-$10) + 1}{}
\end{pfloesung}
\begin{pfloesung}{}
\lz{55.}{ja, $P$ liegt auf beiden Graphen}{f(4) = 16 - 5 = 11, g(4) = 8 + 3 = 11}{}
\end{pfloesung}
\begin{pfloesung}{}
\lz{56.}{$4$}{4; 1; 0; 1}{}
\end{pfloesung}
\begin{pfloesung}{{\small\color{mbgrau}P10 2024 · 3c}}
\lz{57.}{Nein, p($-$5) = 25 $-$ 4 = 21 $\neq$ 20}{($-$5)² = 25}{}
\end{pfloesung}
\begin{pfloesung}{}
\lz{58.}{Anfangswert 45 € , Änderung 38 € je Stunde}{45; 38}{}
\end{pfloesung}
\begin{pfloesung}{}
\lz{59.}{3}{zum Beispiel: 3 : 1 $\Rightarrow$ 3}{}
\end{pfloesung}
\begin{pfloesung}{}
\lz{60.}{$2x + 4$}{y}{}
\end{pfloesung}
\begin{pfloesung}{}
\lz{61.}{Gerade durch $(0|0)$, $(1|3)$, $(2|6)$}{y-Werte: 0, 3, 6}{}
\end{pfloesung}
\begin{pfloesung}{}
\lz{62.}{x Minuten, y Liter}{zum Beispiel ein Tank mit: 300 Litern, der je Minute 15 Liter verliert}{}
\end{pfloesung}
\begin{pfloesung}{}
\lz{63.}{$6x + 10$}{6, n = 10: y}{}
\end{pfloesung}
\begin{pfloesung}{}
\lz{64.}{$48x + 35$}{Nele hat Anfangswert und Änderung vertauscht; richtig: 48, n = 35, also y}{}
\end{pfloesung}
\begin{pfloesung}{{\small\color{mbgrau}P10 2016 · 6c}}
\lz{65.}{y = 1,5x + 200}{}{}
\end{pfloesung}
\begin{pfloesung}{{\small\color{mbgrau}P10 2021 · 7a}}
\lz{66.}{Fahrpreis: y = 0,20x + 5; Fitnesscenter: y = 5x + 20}{20 ct = 0,20 €}{}
\end{pfloesung}
\begin{pfloesung}{{\small\color{mbgrau}P10 2021 · 6c}}
\lz{67.}{y: Höhe in cm; x: Zeit in min; 40: Anfangshöhe in cm}{}{}
\end{pfloesung}
\begin{pfloesung}{{\small\color{mbgrau}P10 2021 · 6d}}
\lz{68.}{200 min (3 h 20 min)}{$-$0,2x + 40 $\Rightarrow$ 0 = 200}{}
\end{pfloesung}
\begin{pfloesung}{{\small\color{mbgrau}P10 2022 · 6b}}
\lz{69.}{y = 55x + 990; 19 Monate}{55x + 990 $\ge$ 2000 $\Rightarrow$ x $\ge$ 18,36; aufrunden}{}
\end{pfloesung}
\begin{pfloesung}{}
\lz{70.}{$0$}{m}{}
\end{pfloesung}
\begin{pfloesung}{}
\lz{71.}{g, denn 5 ist größer als 3: je Schritt nach rechts geht g weiter nach oben}{}{}
\end{pfloesung}
\begin{pfloesung}{}
\lz{72.}{Nein: Beide haben dieselbe Steigung, aber verschiedene y-Achsenabschnitte}{}{}
\end{pfloesung}
\begin{pfloesung}{{\small\color{mbgrau}P10 2021 · 2b}}
\lz{73.}{Aussage 2 und 3 richtig}{f(5) = 16 $\neq$ 0}{}
\end{pfloesung}
\begin{pfloesung}{{\small\color{mbgrau}P10 2023 · 4a}}
\lz{74.}{Gerade durch (0|$-$2) und (1|2); S(1|2); Aussage 1 falsch, Aussage 2 wahr, Aussage 3 falsch}{Steigung: 4, y-Achsenabschnitt $-$2; zweiter Schnittpunkt: ($-$7|$-$30) außerhalb}{}
\end{pfloesung}
\begin{pfloesung}{}
\lz{75.}{Ja: $f(3) = 8$}{}{}
\end{pfloesung}
\begin{pfloesung}{}
\lz{76.}{$-4$}{$\tfrac{3}{4} \cdot (-8) + 2$ $\Rightarrow$ -4}{}
\end{pfloesung}
\begin{pfloesung}{{\small\color{mbgrau}P10 2023 · 1i}}
\lz{77.}{(3|$-$24), denn f(3) = $-$18 $\neq$ $-$24}{f(3) = $-$18}{}
\end{pfloesung}
\begin{pfloesung}{}
\lz{78.}{$265$ €}{$85 + 12 \cdot 15$ $\Rightarrow$ 265}{}
\end{pfloesung}
\begin{pfloesung}{}
\lz{79.}{Nein: $45 + 5 \cdot 52 = 305$ €, das sind 5 € zu viel}{}{}
\end{pfloesung}
\begin{pfloesung}{{\small\color{mbgrau}P10 2021 · 6a}}
\lz{80.}{40 cm; 24 cm}{Abnahme: 2 cm je 10 min}{}
\end{pfloesung}
\begin{pfloesung}{{\small\color{mbgrau}P10 2022 · 6a}}
\lz{81.}{1736 €}{12 · 22 $\Rightarrow$ 264 €}{}
\end{pfloesung}
\begin{pfloesung}{}
\lz{82.}{B}{Startwert auf der y-Achse und Anstieg passen}{}
\end{pfloesung}
\begin{pfloesung}{{\small\color{mbgrau}P10 2023 · 3a}}
\lz{83.}{Angebot 1: C; Angebot 2: B}{Angebot: 1: konstant bis 350 km, dann steigend; Angebot: 2: Grundgebühr, dann gleichmäßig steigend}{}
\end{pfloesung}
\begin{pfloesung}{{\small\color{mbgrau}P10 2016 · 6a}}
\lz{84.}{I Reiselust, II Sonnenschein; II beginnt bei 200 € (Grundpreis)}{200 €}{}
\end{pfloesung}
\begin{pfloesung}{}
\lz{85.}{A ist günstiger}{A: $50 + 12 \cdot 20$ $\Rightarrow$ $290 €, B: 12 \cdot 28$ = 336 €}{}
\end{pfloesung}
\begin{pfloesung}{}
\lz{86.}{das ist nur dann der mit Grundgebühr, wenn er den kleineren Preis je km hat}{Nicht immer: Nicht immer: Rechts vom Schnittpunkt ist der Tarif mit dem kleineren Preis je km günstiger}{}
\end{pfloesung}
\begin{pfloesung}{{\small\color{mbgrau}P10 2016 · 6b}}
\lz{87.}{Reiselust (700 € < 725 €); nein, bei 450 km Sonnenschein 875 € < 900 €}{neuer Wert: 200 + 1,5 · 350 $\Rightarrow$ 725; 2 · 350 $\Rightarrow$ 700}{}
\end{pfloesung}
\begin{pfloesung}{{\small\color{mbgrau}P10 2023 · 3b}}
\lz{88.}{Angebot 1: 178,20 €; Angebot 2: 162,30 €; Angebot 2 günstiger; K(x) = 0,09x + 120}{Prozentwert: 120 Mehrkilometer · 0,36 € $\Rightarrow$ 43,20 €; 5 · 27 $\Rightarrow$ 135 €}{}
\end{pfloesung}
\begin{pfloesung}{\textbf{Prüfstein} \textperiodcentered{} P10 2026 FOR · 5}
\lz{a)}{Gerade durch (0|2) und (1|0), fallend}{Steigungsdreieck: 1 nach rechts, 2 nach unten}{}
\lz{b)}{Ja, f($-$4) = $-$2 · ($-$4) + 2 = 10}{f($-$4) = 10}{}
\lz{c)}{S(2|$-$2)}{2, e = $-$2}{}
\lz{d)}{genau ein gemeinsamer Punkt S(2|$-$2), gleicher Scheitel, p nach oben, q nach unten geöffnet $\Rightarrow$ Berührung nur im Scheitel}{Scheitel beider Parabeln: (2|$-$2)}{}
\end{pfloesung}
\end{document}